%----------------------------------------------------------------------------------------
% Contenido del Informe - Discusión
%----------------------------------------------------------------------------------------

\chapter*{Discusión}
\addcontentsline{toc}{chapter}{Discusión}

El método de ajuste por mínimos cuadrados permitió obtener una función
exponencial que representa la tendencia de los datos experimentales mediante
la minimización de la suma de los cuadrados de los errores entre los valores
medidos y los valores estimados por el modelo.

Debido a que el modelo planteado posee la forma

\[
Q(t)=Ae^{Bt},
\]

fue necesario realizar una transformación logarítmica para convertir la
expresión en un modelo lineal:

\[
\ln(Q)=\ln(A)+Bt.
\]

Esta transformación permitió aplicar el método de mínimos cuadrados lineales
para determinar los coeficientes del modelo. Una vez obtenidos dichos
coeficientes, se recuperó el valor de \(A\) aplicando la función exponencial.

El criterio utilizado por el método consiste en minimizar la suma de los
cuadrados de las diferencias entre los valores observados y los valores
estimados. Al elevar los errores al cuadrado se evita que los errores
positivos y negativos se compensen entre sí y, además, se otorga un mayor
peso a aquellas observaciones que presentan desviaciones más importantes.

La calidad del ajuste puede evaluarse mediante el error cuadrático medio
(ECM), el cual proporciona una medida global de la diferencia entre los datos
experimentales y la función obtenida. Un valor reducido del ECM indica que el
modelo reproduce adecuadamente la tendencia observada en los datos.

En este trabajo, el ajuste exponencial permitió representar el comportamiento
general de la producción del pozo y obtener una función continua a partir de
un conjunto discreto de mediciones. Esto hizo posible realizar estimaciones
para tiempos distintos de los observados experimentalmente y calcular tanto
la fecha estimada de cierre como la producción acumulada y el valor económico
remanente.

Por lo tanto, el método de mínimos cuadrados constituye una herramienta de
gran utilidad para el ajuste de datos experimentales, ya que permite obtener
una función matemática que aproxima el comportamiento real del fenómeno
estudiado y facilita la realización de análisis y predicciones posteriores.
\